% reunion_12_feb.tex

\begin{frame}
    \frametitle{Etapas de diseño de la trayectoria}
    \begin{enumerate}
        \item \textbf{Strategic planning}: Varios días antes.
        \item \textbf{Prediction planning}: Horas antes. Se incluyen las condiciones meteorológicas y posibles conflictos.
        \item \textbf{On route planning}: Durante el vuelo. Detección de conflictos y emergencias.
    \end{enumerate}
\end{frame}


% Frame 1: Introducción a ATFM\
\section{Introducción a ATFM}
\begin{frame}
    \frametitle{Air Traffic Flow Management (ATFM)}
    \subsection{Air Traffic Flow Management (ATFM)}
    \begin{itemize}
        \item Consiste en encontrar horarios y rutas óptimas para \textbf{un conjunto de vuelos} de manera que,
              cuando se ejecuten los planes de vuelo, ninguna región del espacio aéreo tenga más
              aeronaves volando sobre ella que las permitidas por las restricciones de seguridad.
        \item Asimismo, ningún aeropuerto debería tener asignadas más salidas o llegadas de las que puede manejar.
        \item Resuelve el problema de optimizar las salidas y llegadas de un aeropuerto y del espacio aéreo.
        \item Busca garantizar que se cumplan las restricciones de capacidad de salida y
              llegada de cada aeropuerto y que se cumplan las restricciones de capacidad del espacio aéreo.
    \end{itemize}
\end{frame}

% Frame 2: Objetivos de ATFM
\begin{frame}
    \frametitle{Reducción de conflictos}
    \begin{itemize}
        \item Usualmente se busca minimizar el retraso total, en tierra y en aire, imputado a los vuelos.
        \item En este caso, la variable es de tipo binaria.
        \item En \cite{DALSASSO2019} se resuelve un problema multi-objetivo que busca:
              \begin{itemize}
                  \item Minimizar los retrasos.
                  \item Minimizar la desviación de las rutas pre-establecidas.
                  \item Minimizar las tarifas totales de ruta.
              \end{itemize}
              Incluye ponderaciones para ser sometido a las prioridades de las partes interesadas.
    \end{itemize}
\end{frame}

% Frame 3: Problema multi-objetivo en Zhou (2024)
\begin{frame}
    \frametitle{Reducción de conflictos y congestión}
    \begin{itemize}\small
        \item En \cite{Zhou_2024} se resuelve un problema multi-objetivo que busca:
              \begin{itemize}
                  \item Reducir la congestión del espacio aéreo.
                  \item Reducir el consumo de combustible.
                  \item Reducir las posibilidades de conflicto.
              \end{itemize}
        \item Las funciones objetivo a minimizar son:
              \begin{enumerate}
                  \item Minimizar el alcance total horizontal del vuelo (range).
                  \item Minimizar el consumo de combustible.
                  \item Minimizar potenciales conflictos.
                  \item Minimizar la congestión.
                  \item Minimizar la diferencia entre RTA y CTA.
              \end{enumerate}
        \item Para lograr esto, se usan:
              \begin{itemize}
                  \item Variables discretas de posición y altitud.
                  \item Variables continuas de velocidad.
              \end{itemize}
    \end{itemize}
\end{frame}

\begin{frame}
    \begin{figure}
        \centering
        \includegraphics[width=0.9\textwidth]{ng_review_phases}
    \end{figure}



\end{frame}
% Frame 1: Reducción del número de interacciones entre trayectorias
\begin{frame}
    \frametitle{Reducción del número de interacciones entre trayectorias}
    \subsection{Reducción del número de interacciones entre trayectorias}
    \begin{itemize}\small
        \item Como en \cite{Chaimatanan_2014}, se busca minimizar las interacciones de trayectorias de vuelo en todos los puntos de cada trayectoria.
        \item Usualmente se produce en la fase de planificación, y se aplica para las trayectorias en la fase crucero (\textit{En-route}), y ocurre cuando más de una trayectoria compite por el mismo espacio en el mismo periodo de tiempo.
        \item En \cite{Chaimatanan_2014} se hace con datos discretos, y al final se obtiene una ruta óptima y un tiempo de salida (departure) para cada vuelo (30,000 vuelos).
        \item En el caso de \cite{chaimatanan_aircraft_2015}, \cite{Juntama2024} y \cite{alam_metaheuristic_2018}, la metodología propuesta asigna una hora de salida alternativa, una trayectoria de vuelo horizontal y un nivel de vuelo a cada vuelo a escala de país y continente.
    \end{itemize}
\end{frame}

% Frame 2: Estrategias para evitar interacciones
\begin{frame}
    \frametitle{Estrategias para evitar interacciones}
    \begin{table}[h!]
        \centering
        \def\arraystretch{1.5}
        \caption*{\textbf{Estrategias para evitar interacciones}}
        \begin{tabular}{|llc|}
            \hline
            Tiempo de salida alternativo & $t_i = t_{i,0} + \delta_i$                          & \\ \hline
            Trayectoria alternativa      & $w_i = \left(w_i^{1},w_i^{2},\ldots,w_i^{M}\right)$ & \\ \hline
            Nivel de vuelo alternativo   & $FL_i = FL_{i,0} + l_i$                             & \\ \hline
        \end{tabular}
    \end{table}
    \begin{itemize}
        \item \cite{courchelle_simulated_2019} incluye la reducción de conflictos bajo incertidumbre en los parámetros de viento y temperatura.
    \end{itemize}
\end{frame}

% Frame 3: Optimización del perfil vertical y horizontal de las trayectorias
\begin{frame}
    \frametitle{Optimización del perfil vertical y horizontal de las trayectorias}
    Buscan optimizar las trayectorias de referencia en la planificación del vuelo, sobretodo en la etapa crucero. Buscan hacer que el avión vuele a la altura y velocidad óptima, evitando en algunos casos condiciones atmosféricas no favorables.\\
    $$\text{Costo de vuelo} = \text{Combustible} + \text{Cost Index} \times \text{Tiempo de vuelo}$$
    Generalmente se asume que el perfil lateral (LNAV) y vertical (VNAV) se puede modificar durante el vuelo y se adopta el paradigma \textit{``free flight''}, es decir que el piloto puede decidir la ruta sin importar los bloques del espacio aéreo y sus restricciones.
\end{frame}

% Frame 4: Parámetros de configuración y motor
\begin{frame}
    \frametitle{Parámetros de configuración y motor}
    \begin{itemize}
        \item \textbf{Parámetros de configuración de la aeronave}: masa de la aeronave, posición del centro de gravedad, posición del tren de aterrizaje, posición de los flaps/slats, posición de los spoilers/frenos aerodinámicos, etc.;
        \item \textbf{Parámetros del motor}: Estos parámetros son específicos del tipo de motor. Para un motor de reacción, pueden ser el ángulo de palanca de empuje (TLA), velocidades del ventilador del motor, etc.;
        \item \textbf{Condiciones atmosféricas}: temperatura del aire y viento; y
        \item \textbf{Parámetros de la trayectoria de vuelo}: altitud, velocidad, aceleración/desaceleración, ángulo de alabeo, factor de carga, ángulo de ascenso/descenso, tasa de ascenso/descenso, etc.
    \end{itemize}
\end{frame}

% Frame 5: Optimización de trayectorias con WPT
\begin{frame}
    \frametitle{Optimización de trayectorias con WPT}
    \begin{itemize}
        \item En \cite{murrieta-mendoza_particle_2020} se busca encontrar la ruta más económica, tomando en cuenta las restricciones RTA y las condiciones del tiempo. En el caso del tiempo, las condiciones se actualizan en cada WPT a lo largo de la trayectoria.
              \begin{figure}[ht!]
                  \centering
                  \includegraphics[width=0.8\textwidth]{murrieta_WPT_VNAV}
                  \caption*{Waypoints en el perfil vertical (VNAV) \cite{murrieta-mendoza_particle_2020}}
              \end{figure}
    \end{itemize}
\end{frame}

% Frame 6: Malla 2D y prevención de estelas de vapor
\begin{frame}
    \frametitle{Malla 2D y prevención de estelas de vapor}
    \begin{itemize}
        \item En \cite{dancila_botez_2021} se crea una malla en 2D, sobre la que se va a trazar la trayectoria. Además, toma en cuenta las etapas de ascenso y descenso.
              \begin{figure}[ht!]
                  \centering
                  \includegraphics[width=0.6\textwidth]{dancila_grid}
                  \caption*{Ejemplo de una malla \cite{dancila_botez_2021}}
              \end{figure}
        \item Otro enfoque es la prevención de generar estelas de vapor (contrails) a lo largo de la trayectoria.
    \end{itemize}
\end{frame}

% Frame 7: Optimización de trayectoria en condiciones climáticas adversas
\begin{frame}{Optimización de trayectoria en condiciones climáticas adversas}{Sub serf}
    La predicción del mal tiempo y la optimización de trayectorias alternativas se pueden hacer a nivel de uno o varios aviones. En la literatura revisada, se utiliza en gran medida modelos de control predictivo.
    \begin{itemize}
        \item En trabajos como \cite{girardet_generating_2012} y \cite{gonzalez-arribas_robust_2019} se generan trayectorias que evitan condiciones del tiempo de tipo convectivas.
        \item En ambos casos se hace una discretización del espacio. Para las ecuaciones del movimiento y parámetros de los aviones se utilizan datos de BADA.
    \end{itemize}
\end{frame}
\begin{frame}
    \begin{figure}[ht!]
        \centering
        \includegraphics[width=0.6\textwidth]{gonzalez_arribas_storm}
        \caption*{\cite{gonzalez-arribas_robust_2019}}
    \end{figure}




\end{frame}




\begin{frame}
    \scriptsize
    \printbibliography


\end{frame}



